\documentclass[aspectratio=169,11pt]{beamer}
\usepackage[utf8]{inputenc}
\usepackage[T1]{fontenc}
\usepackage[spanish]{babel}
\usepackage[sfdefault]{FiraSans}
\usepackage[scale=0.9]{FiraMono}
\usepackage{fontawesome5}
\usepackage{tikz}
\usetikzlibrary{arrows.meta,positioning,fit,calc,shadings}
\usepackage{booktabs,array,listings,colortbl}
\usepackage{pgfplots}
\pgfplotsset{compat=1.18}
\usetheme{metropolis}
\metroset{block=fill,numbering=fraction,progressbar=frametitle}

% ------------------------------------------------------------ paleta oscuro tech
\definecolor{fondo}{HTML}{14131C}
\definecolor{panel}{HTML}{211F2E}
\definecolor{panel2}{HTML}{2D2A40}
\definecolor{texto}{HTML}{ECEAF5}
\definecolor{tenue}{HTML}{A9A5C0}
\definecolor{vl}{HTML}{A59EFF}
\definecolor{menta}{HTML}{5EE6D0}
\definecolor{ambar}{HTML}{FFC266}
\definecolor{rojo}{HTML}{FF7A85}
\definecolor{verde}{HTML}{7EE29A}

\setbeamercolor{background canvas}{bg=fondo}
\setbeamercolor{normal text}{fg=texto,bg=fondo}
\setbeamercolor{structure}{fg=vl}
\setbeamercolor{frametitle}{fg=texto,bg=panel}
\setbeamercolor{progress bar}{fg=menta,bg=panel2}
\setbeamercolor{progress bar in head/foot}{fg=menta,bg=panel2}
\setbeamercolor{alerted text}{fg=ambar}
\setbeamercolor{block title}{fg=vl,bg=panel2}
\setbeamercolor{block body}{fg=texto,bg=panel}
\setbeamercolor{block title alerted}{fg=ambar,bg=panel2}
\setbeamercolor{block body alerted}{fg=texto,bg=panel}
\setbeamercolor{itemize item}{fg=menta}
\setbeamercolor{itemize subitem}{fg=vl}
\setbeamercolor{enumerate item}{fg=vl}
\setbeamercolor{description item}{fg=vl}
\setbeamercolor{footline}{fg=tenue}
\setbeamercolor{page number in head/foot}{fg=tenue}
\setbeamerfont{block title}{series=\bfseries}
\setbeamertemplate{frame footer}{\faNetworkWired\enspace Equipo 5 \textperiodcentered{} Memoria compartida distribuida}
\arrayrulecolor{tenue!45!fondo}

% fondo: cuadrícula de puntos muy tenue
\setbeamertemplate{background}{%
  \begin{tikzpicture}
    \useasboundingbox (0,0) rectangle (\paperwidth,\paperheight);
    \foreach \x in {0.3,0.8,...,16}{\foreach \y in {0.3,0.8,...,9}{\fill[vl,opacity=.09] (\x,\y) circle (.4pt);}}
  \end{tikzpicture}}

\lstset{basicstyle=\ttfamily\scriptsize\color{texto},columns=fullflexible,
  keywordstyle=\bfseries\color{vl},commentstyle=\color{tenue}\itshape,
  emph={LEER,ESCRIBIR,INVALIDA,ACK,DATOS,R,W},emphstyle=\bfseries\color{menta},
  morekeywords={si,return,para,en,sino,al,recibir,de},morecomment=[l]{//},
  backgroundcolor=\color{panel},frame=leftline,framerule=2.5pt,rulecolor=\color{vl},
  xleftmargin=.7em,framexleftmargin=.5em,aboveskip=.4em,belowskip=.4em}

\newcolumntype{L}[1]{>{\raggedright\arraybackslash}p{#1}}

% ------------------------------------------------------------ piezas gráficas
% red decorativa de laptops
\newcommand{\redDeco}[1]{%
  \begin{scope}[#1]
    \coordinate (n1) at (0,0);     \coordinate (n2) at (1.9,1.15);  \coordinate (n3) at (3.4,.1);
    \coordinate (n4) at (1.4,-1.35);\coordinate (n5) at (3.1,-2.0);  \coordinate (n6) at (4.7,-1.0);
    \coordinate (n7) at (.1,-2.7); \coordinate (n8) at (4.3,1.8);   \coordinate (n9) at (5.5,.45);
    \coordinate (n10) at (1.9,-3.5);\coordinate (n11) at (5.0,-3.0); \coordinate (n12) at (.6,1.9);
    \foreach \a/\b in {n1/n2,n2/n3,n1/n4,n4/n3,n3/n6,n4/n5,n5/n6,n2/n8,n8/n9,n3/n9,n6/n9,n4/n7,
                       n7/n10,n5/n10,n5/n11,n6/n11,n3/n8,n12/n2,n12/n1}
      {\draw[vl!55!fondo,line width=.9pt] (\a) -- (\b);}
    \foreach \a/\b/\t in {n1/n2/.45,n3/n6/.6,n4/n5/.35,n8/n9/.5,n5/n11/.55}
      {\fill[menta] ($(\a)!\t!(\b)$) circle (1.5pt);}
    \foreach \n in {n2,n5,n7,n9,n10,n11,n12}
      {\fill[vl,opacity=.25] (\n) circle (6pt); \fill[vl] (\n) circle (2.4pt);}
    \foreach \n/\ic/\c in {n1/\faLaptop/menta,n4/\faLaptop/menta,n6/\faLaptop/menta,n8/\faLaptop/menta,n3/\faServer/ambar}
      {\fill[\c,opacity=.14] (\n) circle (.62);
       \node[circle,draw=\c,line width=1.1pt,fill=panel,text=\c,minimum size=.8cm,inner sep=0pt] at (\n) {\small\ic};}
  \end{scope}}

% barra de nivel (0-5)
\newcommand{\nivel}[2]{\tikz[baseline=-.1ex]{\foreach \k in {1,...,5}{%
  \ifnum\k>#1\relax\fill[panel2] (\k*.2,0) rectangle +(.15,.22);\else\fill[#2] (\k*.2,0) rectangle +(.15,.22);\fi}}}

\tikzset{node/.style={draw=#1,fill=panel,line width=1pt,rounded corners=4pt,minimum width=3.4cm,minimum height=1.05cm,align=center,text=texto},
  node/.default=vl,
  box/.style={draw=#1,fill=panel,line width=1pt,rounded corners=4pt,align=center,minimum width=3.1cm,minimum height=.95cm,font=\footnotesize,text=texto}}
\tikzset{stat/.style={rounded corners=6pt,fill=panel,draw=#1,line width=1.2pt,align=center,inner sep=7pt}}

% secciones con expositor
\newcommand{\exponente}{}
\newcommand{\seccion}[2]{\renewcommand{\exponente}{#2}\section{#1}}

\setbeamertemplate{section page}{%
  \begin{tikzpicture}[remember picture,overlay]
    \redDeco{shift={($(current page.north east)+(-6.1,-3.3)$)},opacity=.6}
    \node[anchor=south west,inner sep=0pt,font=\fontsize{66}{66}\selectfont\bfseries,text=vl]
      (num) at ($(current page.west)+(1.1,.35)$) {0\insertsectionnumber};
    \draw[menta,line width=2pt] ($(num.south west)+(0,-.25)$) -- ++(2.0,0);
    \node[anchor=north west,inner sep=0pt,text width=8.2cm,font=\LARGE\bfseries,text=texto]
      (tit) at ($(num.south west)+(0,-.55)$) {\insertsectionhead};
    \node[anchor=north west,inner sep=0pt,text=ambar] at ($(tit.south west)+(0,-.35)$) {\faUser\enspace Expone: \exponente};
  \end{tikzpicture}}

\title{Memoria compartida distribuida y modelos de consistencia}
\subtitle{Cómputo paralelo -- Tema 5}
\author[Equipo 5]{Ever Garcia Ramirez, Christian Luna Ortiz, Marcelino Parra Hernández, José José Varela Palacios}
\date{10 de noviembre de 2026}

\begin{document}

% ============================================================ PORTADA
\begin{frame}[plain]
  \begin{tikzpicture}[remember picture,overlay]
    \fill[panel] (current page.south west) rectangle ($(current page.south east)+(0,1.75)$);
    \fill[vl] ($(current page.south west)+(0,1.75)$) rectangle ($(current page.south east)+(0,1.79)$);
    \redDeco{shift={($(current page.north east)+(-5.75,-2.75)$)},scale=.93}
    \node[anchor=north west,text width=9.4cm,inner sep=0pt] at ($(current page.north west)+(1,-1.15)$) {%
      {\color{menta}\footnotesize\bfseries\faNetworkWired\enspace CÓMPUTO PARALELO Y DISTRIBUIDO \textperiodcentered{} TEMA 5}\par\vspace{.55em}
      {\fontsize{25}{29}\selectfont\bfseries\color{texto}Memoria compartida\\distribuida\par}\vspace{.3em}
      {\Large\color{vl}y modelos de consistencia\par}\vspace{.7em}
      {\color{vl}\rule{2.2cm}{2pt}}\par\vspace{.6em}
      {\small\color{tenue}\faCalendar\enspace 10 de noviembre de 2026\qquad\faUsers\enspace Equipo 5}};
    \node[anchor=west,text=texto,inner sep=0pt] at ($(current page.south west)+(1,.88)$) {%
      \footnotesize\begin{tabular}{@{}l@{\hspace{2.6em}}l@{}}
        {\color{menta}\faUser}\enspace Ever Garcia Ramirez & {\color{menta}\faUser}\enspace Christian Luna Ortiz\\[.25em]
        {\color{menta}\faUser}\enspace Marcelino Parra Hernández & {\color{menta}\faUser}\enspace José José Varela Palacios
      \end{tabular}};
  \end{tikzpicture}
\end{frame}

% ============================================================ CONTENIDO
\begin{frame}{Contenido}
  \centering
  \begin{tikzpicture}
    \foreach \i/\t/\who in {1/{Qué es la DSM y cómo se simula}/Ever,
                            2/{Esquemas: páginas, variables y objetos}/Christian,
                            3/{Modelos de consistencia}/Marcelino,
                            4/{Acceso del Equipo A a datos del Equipo B}/José,
                            5/{Prueba real en un clúster de laptops}/Christian}{
      \node[fill=panel,rounded corners=5pt,minimum width=12.8cm,minimum height=.92cm] at (0,-\i*1.08) {};
      \fill[vl] (-6.4,-\i*1.08-.46) rectangle +(.07,.92);
      \node[font=\large\bfseries,text=vl] at (-5.75,-\i*1.08) {0\i};
      \node[anchor=west,text=texto] at (-5.15,-\i*1.08) {\t};
      \node[anchor=east,text=ambar,font=\footnotesize] at (6.15,-\i*1.08) {\faUser\enspace\who};
    }
  \end{tikzpicture}
\end{frame}

%%%%%%%%%%%%%%%%%%%%%%%%%%%%%%%%%%%%%%%%%%%%%%%%%%%%%%%%%%%%%%%%%%
\seccion{Qué es la DSM y cómo se simula}{Ever}

\begin{frame}{Dos formas de comunicar procesos}
  \begin{columns}[T]
    \column{.48\textwidth}
    \begin{block}{\faMemory\enspace Memoria compartida}
      \begin{itemize}
        \item Procesos leen y escriben en un mismo espacio de direcciones.
        \item Programación simple.
        \item Requiere hardware común (multiprocesador).
      \end{itemize}
    \end{block}
    \column{.48\textwidth}
    \begin{block}{\faEnvelope\enspace Paso de mensajes (MPI)}
      \begin{itemize}
        \item Cada nodo tiene memoria propia.
        \item Comunicación explícita: \texttt{send}/\texttt{recv}.
        \item Escala en clústeres de computadoras separadas.
      \end{itemize}
    \end{block}
  \end{columns}
  \vspace{.9em}
  {\color{ambar}\faLightbulb}\enspace\alert{Pregunta:} ¿podemos tener la comodidad de la primera sobre el hardware de la segunda?
\end{frame}

\begin{frame}{Memoria compartida distribuida (DSM)}
  \begin{itemize}
    \item \textbf{Abstracción}: una capa de software que presenta a los procesos de varias computadoras \emph{un solo espacio de direcciones} compartido.
    \item Por debajo, los datos viven en la memoria local de cada nodo y viajan por la red cuando se necesitan.
  \end{itemize}
  \begin{block}{\faCogs\enspace Cómo se simula}
    \begin{enumerate}
      \item Se divide el espacio compartido en unidades (páginas, variables u objetos).
      \item Se registra qué nodo tiene cada unidad (directorio).
      \item Un acceso a una unidad remota se intercepta y se resuelve por la red.
      \item Se guardan copias locales y se mantienen coherentes.
    \end{enumerate}
  \end{block}
\end{frame}

\begin{frame}{Problemas que hay que resolver}
  \begin{itemize}\setlength{\itemsep}{.55em}
    \item[\textcolor{ambar}{\faClock}] \textbf{Latencia}: acceder a un dato remoto cuesta órdenes de magnitud más que uno local.
    \item[\textcolor{ambar}{\faSearch}] \textbf{Localización}: ¿quién tiene la copia actual? (directorio central o distribuido).
    \item[\textcolor{ambar}{\faCopy}] \textbf{Replicación}: copias múltiples aceleran lecturas, pero complican las escrituras.
    \item[\textcolor{ambar}{\faIcon{sync-alt}}] \textbf{Coherencia y consistencia}: qué valor ve cada proceso y en qué orden.
    \item[\textcolor{ambar}{\faRandom}] \textbf{Falso compartimiento}: dos variables distintas en la misma unidad provocan tráfico innecesario.
  \end{itemize}
\end{frame}

%%%%%%%%%%%%%%%%%%%%%%%%%%%%%%%%%%%%%%%%%%%%%%%%%%%%%%%%%%%%%%%%%%
\seccion{Esquemas: páginas, variables y objetos}{Christian}

\begin{frame}{Tres granularidades}
  \centering\small
  \begin{tabular}{@{}>{\color{menta}\bfseries}L{2.1cm}L{3.5cm}L{3.5cm}L{3.5cm}@{}}
    \toprule
    & {\color{vl}\faIcon{file-alt}}\enspace\textbf{Páginas} & {\color{vl}\faTag}\enspace\textbf{Variables} & {\color{vl}\faCubes}\enspace\textbf{Objetos} \\ \midrule
    Unidad & Página de memoria virtual (p.~ej. 4~KB) & Variable compartida declarada & Objeto con operaciones \\[.25em]
    Detección & Falla de página (MMU) & Compilador o biblioteca & Invocación de método \\[.25em]
    Ejemplos & IVY, TreadMarks & Munin & Orca, Linda (tuplas) \\[.25em]
    Ventaja & Transparente, sin cambiar el programa & Grano fino, menos falso compartimiento & Encapsula datos y coherencia \\[.25em]
    Desventaja & Falso compartimiento & Hay que anotar el código & Modelo de programación nuevo \\
    \bottomrule
  \end{tabular}
\end{frame}

\begin{frame}{DSM basada en páginas}
  \begin{itemize}
    \item Se aprovecha la MMU: las páginas remotas se marcan como \emph{no presentes} o de \emph{solo lectura}.
    \item Una falla de página dispara el protocolo de DSM, que trae la página y reintenta el acceso.
    \item Un \emph{dueño} (o gestor) conoce el estado de cada página: lectura compartida o escritura exclusiva.
  \end{itemize}
  \begin{block}{\faIcon{shield-alt}\enspace Coherencia típica}
    \begin{itemize}
      \item \textbf{Escritura-invalidación}: antes de escribir se invalidan las demás copias.
      \item \textbf{Escritura-actualización}: la escritura se propaga a todas las copias.
    \end{itemize}
  \end{block}
\end{frame}

\begin{frame}{DSM basada en variables y en objetos}
  \begin{columns}[T]
    \column{.48\textwidth}
    \begin{block}{\faTag\enspace Variables compartidas (Munin)}
      \begin{itemize}
        \item Solo se comparte lo declarado; cada variable usa el protocolo adecuado a su patrón (solo lectura, migratoria, productor-consumidor\ldots).
        \item Cerrojos y barreras: los cambios se propagan al liberar.
      \end{itemize}
    \end{block}
    \column{.48\textwidth}
    \begin{block}{\faCubes\enspace Objetos (Orca, Linda)}
      \begin{itemize}
        \item Los datos solo se tocan mediante operaciones; el sistema replica o migra el objeto.
        \item Operaciones atómicas: consistencia más simple.
        \item Linda: se intercambian \emph{tuplas} con \texttt{out}, \texttt{in}, \texttt{rd}.
      \end{itemize}
    \end{block}
  \end{columns}
\end{frame}

%%%%%%%%%%%%%%%%%%%%%%%%%%%%%%%%%%%%%%%%%%%%%%%%%%%%%%%%%%%%%%%%%%
\seccion{Modelos de consistencia}{Marcelino}

\begin{frame}{Coherencia vs.\ consistencia}
  \begin{itemize}
    \item \textbf{Coherencia}: qué valores puede devolver una lectura de \emph{una misma} ubicación.
    \item \textbf{Consistencia}: en qué orden se ven las operaciones sobre \emph{distintas} ubicaciones.
  \end{itemize}
  \begin{alertblock}{\faHandshake\enspace Un modelo de consistencia es un contrato}
    \begin{itemize}
      \item Más estricto $\Rightarrow$ más fácil de programar, pero más lento (más mensajes y esperas).
      \item Más relajado $\Rightarrow$ más rendimiento, pero el programador sincroniza explícitamente.
    \end{itemize}
  \end{alertblock}
\end{frame}

\begin{frame}{Modelos de consistencia}
  \begin{columns}[c]
    \column{.07\textwidth}
    \begin{tikzpicture}
      \shade[top color=vl,bottom color=menta] (0,0) rectangle (.16,-4.6);
      \fill[menta] (-.12,-4.6) -- (.28,-4.6) -- (.08,-4.95) -- cycle;
      \node[above,font=\tiny\bfseries,text=vl,align=center] at (.08,0) {más\\estricto};
      \node[below,font=\tiny\bfseries,text=menta,align=center] at (.08,-4.95) {más\\relajado};
    \end{tikzpicture}
    \column{.9\textwidth}
    \small
    \begin{description}\setlength{\itemsep}{.35em}
      \item[Estricta] Toda lectura devuelve la última escritura en tiempo real. Imposible de lograr en una red real.
      \item[Secuencial (Lamport)] El resultado es el de algún orden secuencial único que respeta el orden de cada proceso; todos ven el mismo orden.
      \item[Causal] Solo se garantiza el orden de operaciones relacionadas causalmente; las concurrentes pueden verse en distinto orden.
      \item[Débil / de liberación] El orden solo se garantiza en puntos de sincronización (\texttt{acquire}/\texttt{release}).
      \item[Eventual] Si dejan de haber escrituras, todas las réplicas convergen con el tiempo.
    \end{description}
  \end{columns}
\end{frame}

\begin{frame}{Ejemplo: secuencial vs.\ causal}
  \begin{columns}[T]
    \column{.5\textwidth}
    \begin{block}{\faCode\enspace Programa}
      \texttt{P1: x = 1}\\
      \texttt{P2: if (x == 1) y = 1}\\
      \texttt{P3: lee y, luego x}
    \end{block}
    \textbf{Pregunta:} si P3 lee \texttt{y == 1}, ¿puede leer después \texttt{x == 0}?
    \column{.5\textwidth}
    \begin{itemize}\setlength{\itemsep}{.5em}
      \item[\textcolor{verde}{\faTimesCircle}] \alert{Secuencial}: no. Todos ven el mismo orden.
      \item[\textcolor{verde}{\faTimesCircle}] \alert{Causal}: no. \texttt{y=1} depende causalmente de \texttt{x=1}.
      \item[\textcolor{rojo}{\faCheckCircle}] \alert{Eventual}: sí, momentáneamente.
    \end{itemize}
  \end{columns}
\end{frame}

\begin{frame}{Comparación}
  \centering
  \begin{tabular}{@{}>{\bfseries}l l l l@{}}
    \toprule
    \textbf{Modelo} & \textbf{Facilidad} & \textbf{Rendimiento} & \textbf{Sincronización} \\ \midrule
    Estricta           & \nivel{5}{vl}\enspace Máxima & \nivel{1}{menta}\enspace Muy bajo & No necesaria \\[.2em]
    Secuencial         & \nivel{4}{vl}\enspace Alta   & \nivel{2}{menta}\enspace Bajo     & No necesaria \\[.2em]
    Causal             & \nivel{3}{vl}\enspace Media  & \nivel{3}{menta}\enspace Medio    & Parcial \\[.2em]
    Débil / liberación & \nivel{2}{vl}\enspace Baja   & \nivel{4}{menta}\enspace Alto     & Explícita \\[.2em]
    Eventual           & \nivel{2}{vl}\enspace Baja   & \nivel{5}{menta}\enspace Muy alto & Lógica de aplicación \\
    \bottomrule
  \end{tabular}
\end{frame}

%%%%%%%%%%%%%%%%%%%%%%%%%%%%%%%%%%%%%%%%%%%%%%%%%%%%%%%%%%%%%%%%%%
\seccion{Acceso del Equipo A a datos del Equipo B}{José}

\begin{frame}{Cómo accede el Equipo A a datos del Equipo B}
  \centering
  \begin{tikzpicture}[>={Stealth[length=2.4mm]}]
    \node[node] (pa) at (0,2.3) {{\color{vl}\faCode}\enspace Proceso en A\\\texttt{x = lee(dir)}};
    \node[node=menta] (dsm) at (0,0.5) {Capa DSM de A\\{\footnotesize\color{tenue}(directorio, caché)}};
    \node[node=menta] (dsmb) at (8,0.5) {Capa DSM de B\\{\footnotesize\color{tenue}(directorio)}};
    \node[node=verde] (mem) at (8,2.3) {{\color{verde}\faMemory}\enspace Memoria local de B\\{\footnotesize\color{tenue}(dato en \texttt{dir})}};
    \node[draw=tenue!60!fondo,dashed,rounded corners=6pt,fit=(pa)(dsm),inner sep=7pt,
          label={[text=menta,font=\footnotesize\bfseries]above:\faLaptop\ Nodo A}] {};
    \node[draw=tenue!60!fondo,dashed,rounded corners=6pt,fit=(dsmb)(mem),inner sep=7pt,
          label={[text=menta,font=\footnotesize\bfseries]above:\faLaptop\ Nodo B}] {};
    \draw[->,tenue,line width=.9pt] (pa) -- node[right,font=\scriptsize]{1. acceso} (dsm);
    \draw[->,menta,line width=1.1pt] ([yshift=4pt]dsm.east) -- node[above,font=\scriptsize]{2. solicitud por red} ([yshift=4pt]dsmb.west);
    \draw[->,tenue,line width=.9pt] (dsmb) -- node[left,font=\scriptsize]{3. lee} (mem);
    \draw[->,verde,line width=1.1pt] ([yshift=-4pt]dsmb.west) -- node[below,font=\scriptsize]{4. respuesta con el dato} ([yshift=-4pt]dsm.east);
    \draw[->,dashed,verde,line width=.9pt] (dsm.north west) to[bend left=30] node[left,font=\scriptsize]{5. entrega} (pa.south west);
    \node[font=\footnotesize,text=tenue] at (4,-1.05) {\faInfoCircle\enspace Para el proceso en A todo parece una lectura local de memoria.};
  \end{tikzpicture}
\end{frame}

\begin{frame}[fragile]{Pseudocódigo: lo que hace cada nodo}
\begin{columns}[T]
\column{.49\textwidth}
\begin{lstlisting}
leer(dir):
  p = pagina(dir)
  si p esta en cache:             // R o W
      return cache[p][dir]
  enviar(gestor, LEER, p)         // falla
  datos = recibir(gestor)
  cache[p] = (R, datos)           // copia
  return cache[p][dir]
\end{lstlisting}
\column{.49\textwidth}
\begin{lstlisting}
escribir(dir, v):
  p = pagina(dir)
  si cache[p].modo != W:
      // pedir propiedad EXCLUSIVA
      enviar(gestor, ESCRIBIR, p)
      datos = recibir(gestor)
      cache[p] = (W, datos)
  cache[p][dir] = v
\end{lstlisting}
\end{columns}
\vspace{.4em}
{\small{\color{ambar}\faLock}\enspace El nodo nunca escribe sin ser el único dueño: así dos escrituras simultáneas no se pisan.}
\end{frame}

\begin{frame}[fragile]{Pseudocódigo: el gestor (escritura-invalidación)}
\begin{lstlisting}
al recibir ESCRIBIR(p) de n:
  si duenio[p] != ninguno y duenio[p] != n:
      datos[p] = traer_e_invalidar(duenio[p], p)   // ultima version
  para c en copias[p] - {n}:
      enviar(c, INVALIDA, p);  esperar ACK de c     // nadie lee valores viejos
  duenio[p] = n;  copias[p] = {}
  responder(n, DATOS, datos[p], modo = W)

al recibir LEER(p) de n:
  si duenio[p] != ninguno:
      datos[p] = traer(duenio[p], p)                // el dueno baja a modo R
  copias[p] = copias[p] + {n}
  responder(n, DATOS, datos[p], modo = R)
\end{lstlisting}
\end{frame}

%%%%%%%%%%%%%%%%%%%%%%%%%%%%%%%%%%%%%%%%%%%%%%%%%%%%%%%%%%%%%%%%%%
\seccion{Prueba real en un clúster de laptops}{Christian}

\begin{frame}{Nuestro clúster: maestro--trabajador}
  \centering
  \begin{tikzpicture}[>={Stealth[length=2.2mm]}]
    \node[box=ambar] (m) at (0,2.5) {{\color{ambar}\faLaptop}\enspace\textbf{Laptop maestra}\\gestor DSM: directorio + páginas\\{\color{ambar}\texttt{192.168.1.10:5000}}};
    \node[draw=vl,line width=1pt,circle,fill=panel,text=vl,minimum size=.95cm] (sw) at (0,.85) {\faWifi};
    \node[box=menta] (w1) at (-4.6,-.65) {{\color{menta}\faLaptop}\enspace Trabajador 1\\{\color{menta}\faMicrochip\ CPU}\enspace\texttt{192.168.1.11}};
    \node[box=verde] (w2) at (0,-.65) {{\color{verde}\faLaptop}\enspace Trabajador 2\\{\color{verde}\faBolt\ GPU}\enspace\texttt{192.168.1.12}};
    \node[box=verde] (w3) at (4.6,-.65) {{\color{verde}\faLaptop}\enspace Trabajador 3\\{\color{verde}\faBolt\ GPU}\enspace\texttt{192.168.1.13}};
    \draw[<->,menta,line width=1pt] (m) -- (sw);
    \foreach \w in {w1,w2,w3} \draw[<->,menta,line width=1pt] (sw) -- (\w.north);
    \node[font=\scriptsize,align=left,anchor=west,text=tenue] at (2.9,1.2) {{\color{menta}\faIcon{exchange-alt}}\ Sockets TCP:\\\texttt{LEER}, \texttt{ESCRIBIR},\\\texttt{INVALIDA}, \texttt{ACK}};
  \end{tikzpicture}\par
  \vspace{.4em}
  {\footnotesize\color{tenue}Prototipo propio en Python (solo biblioteca estándar): DSM por páginas de 4~KB, gestor central y escritura-invalidación. IPs ilustrativas; las mediciones de las siguientes diapositivas se hicieron con este mismo programa.}
\end{frame}

\begin{frame}{Latencia: el costo de una falla de página}
  \begin{block}{\faCalculator\enspace Modelo}
    \[
      T_{\text{falla}} \;\approx\; T_{\text{RTT}} + \frac{S_{\text{página}}}{B} + T_{\text{gestor}}
      \qquad\qquad
      \text{penalización} = \frac{T_{\text{falla}}}{T_{\text{local}}}
    \]
  \end{block}
  \centering
  \begin{tikzpicture}
    \node[stat=menta,minimum width=6.3cm,minimum height=2.35cm] at (-3.45,0) {%
      {\footnotesize\color{tenue}\faHome\enspace Misma máquina (\emph{loopback})}\\[4pt]
      {\fontsize{28}{30}\selectfont\bfseries\color{menta}$\approx$\,700$\times$}\\[5pt]
      {\scriptsize RTT 0.06--0.10 ms\quad\textperiodcentered\quad falla 0.21 ms}};
    \node[stat=rojo,minimum width=6.3cm,minimum height=2.35cm] at (3.45,0) {%
      {\footnotesize\color{tenue}\faGlobeAmericas\enspace Internet (laptop $\leftrightarrow$ servidor, Tailscale)}\\[4pt]
      {\fontsize{28}{30}\selectfont\bfseries\color{rojo}$\approx$\,376\,000$\times$}\\[5pt]
      {\scriptsize RTT 111 ms\quad\textperiodcentered\quad falla 113 ms}};
  \end{tikzpicture}\\[.45em]
  {\footnotesize\color{tenue}Veces más lento que una lectura local ($T_{\text{local}} \approx 0.30\,\mu$s, página ya en caché, en Python). La falla casi es solo $T_{\text{RTT}}$: \textcolor{ambar}{la red domina}.}
\end{frame}

\begin{frame}{Falso compartimiento, medido}
  \begin{columns}[c]
    \column{.55\textwidth}
    \begin{tikzpicture}
      \begin{axis}[ybar, ymode=log, log origin=infty, width=7.2cm, height=5.4cm,
          symbolic x coords={Loopback,Internet}, xtick=data,
          ylabel={incrementos/s (2 nodos)}, ylabel style={font=\scriptsize,text=tenue},
          ymin=1e4, ymax=1e7, bar width=16pt, enlarge x limits=0.45,
          axis line style={tenue}, tick style={tenue}, ymajorgrids, major grid style={panel2},
          axis x line*=bottom, axis y line*=left,
          legend style={font=\scriptsize,at={(.5,-.2)},anchor=north,legend columns=2,draw=none,fill=none,text=texto},
          tick label style={font=\scriptsize,text=texto},
          point meta=explicit symbolic, nodes near coords,
          every node near coord/.append style={font=\scriptsize,text=texto}]
        \addplot[fill=verde,draw=none] coordinates {(Loopback,2531000)[2.53\,M] (Internet,2520000)[2.52\,M]};
        \addplot[fill=rojo,draw=none] coordinates {(Loopback,1143000)[1.14\,M] (Internet,33300)[33\,k]};
        \legend{Páginas distintas, Misma página}
      \end{axis}
    \end{tikzpicture}
    \column{.43\textwidth}
    \footnotesize
    Dos trabajadores incrementan \emph{cada uno su propia variable} al mismo tiempo.
    \begin{itemize}
      \item[\textcolor{verde}{\faCheckCircle}] Páginas distintas: $0$ fallas.
      \item[\textcolor{rojo}{\faRandom}] Misma página: la página rebota entre nodos (invalidación tras invalidación).
    \end{itemize}
    \vspace{.3em}
    \begin{tikzpicture}
      \node[stat=ambar,minimum width=2.45cm,minimum height=1.45cm] at (0,0) {{\scriptsize\color{tenue}loopback}\\{\LARGE\bfseries\color{ambar}2.2$\times$}};
      \node[stat=rojo,minimum width=2.45cm,minimum height=1.45cm] at (2.75,0) {{\scriptsize\color{tenue}internet}\\{\LARGE\bfseries\color{rojo}76$\times$}};
    \end{tikzpicture}\\[.3em]
    {\color{tenue}más lento con la misma página.} Resultado final correcto en todas las corridas.
  \end{columns}
\end{frame}

\begin{frame}{Balanceo CPU/GPU con DSM}
  \begin{columns}[T]
    \column{.48\textwidth}
    \begin{block}{\faBalanceScale\enspace El dueño calcula (\emph{owner computes})}
      \begin{itemize}
        \item El maestro reparte \emph{páginas}, no tareas sueltas: cada nodo calcula sobre lo que ya tiene.
        \item Laptops con GPU reciben bloques más grandes, en proporción a su velocidad medida.
      \end{itemize}
    \end{block}
    \column{.48\textwidth}
    \begin{block}{\faClipboardCheck\enspace Reglas que salen de la medición}
      \begin{itemize}
        \item Alinear los datos de cada nodo a \emph{su} página: evita el $76\times$.
        \item La VRAM no forma parte de la DSM: el nodo copia su página a la GPU, calcula y escribe de vuelta.
        \item Consistencia de liberación: enviar cambios en lotes al \texttt{release}, no en cada escritura.
      \end{itemize}
    \end{block}
  \end{columns}
\end{frame}

\begin{frame}{Conclusiones}
  \begin{itemize}\setlength{\itemsep}{.45em}
    \item[\textcolor{menta}{\faCheck}] La DSM da la comodidad de memoria compartida sobre computadoras con memoria separada.
    \item[\textcolor{menta}{\faCheck}] La elección de granularidad (páginas, variables, objetos) equilibra transparencia, falso compartimiento y esfuerzo de programación.
    \item[\textcolor{menta}{\faCheck}] El modelo de consistencia es el compromiso central entre facilidad de uso y rendimiento.
    \item[\textcolor{menta}{\faCheck}] Lo medimos: en red, una falla cuesta $\sim\!10^5$ veces una lectura local; acomodar los datos importa más que el procesador.
    \item[\textcolor{menta}{\faCheck}] Por eso en clústeres reales domina el paso de mensajes (MPI): el costo de la red queda explícito.
  \end{itemize}
\end{frame}

\begin{frame}{Referencias}
  \footnotesize
  \begin{itemize}\setlength{\itemsep}{.4em}
    \item[\textcolor{vl}{\faBook}] Tanenbaum, A. y Van Steen, M. \emph{Sistemas distribuidos: principios y paradigmas}.
    \item[\textcolor{vl}{\faBook}] Li, K. y Hudak, P. (1989). Memory coherence in shared virtual memory systems. \emph{ACM TOCS}.
    \item[\textcolor{vl}{\faBook}] Lamport, L. (1979). How to make a multiprocessor computer that correctly executes multiprocess programs. \emph{IEEE Trans.\ Computers}.
    \item[\textcolor{vl}{\faBook}] Bennett, J. et al.\ (1990). Munin: distributed shared memory based on type-specific memory coherence. \emph{PPOPP}.
    \item[\textcolor{vl}{\faBook}] Keleher, P. et al.\ (1994). TreadMarks: distributed shared memory on standard workstations and operating systems. \emph{USENIX}.
  \end{itemize}
\end{frame}

\begin{frame}[plain]
  \begin{tikzpicture}[remember picture,overlay]
    \redDeco{shift={($(current page.north east)+(-5.75,-2.75)$)},scale=.93}
    \node[anchor=west,inner sep=0pt] at ($(current page.west)+(1.1,.6)$) {%
      \begin{minipage}{8cm}
        {\fontsize{40}{44}\selectfont\bfseries\color{texto}¿Preguntas?\par}\vspace{.4em}
        {\color{vl}\rule{2.2cm}{2pt}}\par\vspace{.6em}
        {\large\color{menta}Gracias\par}\vspace{.4em}
        {\small\color{tenue}Equipo 5 \textperiodcentered{} Memoria compartida distribuida}
      \end{minipage}};
  \end{tikzpicture}
\end{frame}

\end{document}
